% Unregistered strengthening of NP; not an ordinary-Chow survival assertion.
\paragraph{The elliptic periods are outside the component Picard span
(audited research lemma).}
Retain all notation of (NP.1)--(NP.8). Let $C=C_{12}$, and let
$r_1,\ldots,r_{24}$ be the zero divisor on $C$ of the generic sextic $G$.
Then, for $j=2,3,4$,
\begin{equation}
 [q_j-q_1]\notin\operatorname{im}\bigl(
 \operatorname{Pic}(E_1)_{\mathbf Q}\oplus
 \operatorname{Pic}(E_2)_{\mathbf Q}\longrightarrow
 \operatorname{Pic}(C)_{\mathbf Q}\bigr).
                                                        \tag{NP.9}
\end{equation}
This is a statement about component divisor restrictions. It does
not identify the image of all proper perfect-complex families and
does not assert that a class survives in the ordinary Chow quotient.

\emph{Proof.}
Work first over $K=\overline{k_0(t_1,t_2,t_3)}$, so that the four
points $q_j$ and the marked conic intersections are already defined.
The restriction map
\[
 H^0(\mathbf P^3_K,\mathcal O(6))\longrightarrow
 H^0(C_K,\mathcal O_C(6))
\]
is onto. Indeed the Koszul resolution of the complete intersection
of two quadrics gives
$0\to\mathcal O(2)\to\mathcal O(4)^{\oplus2}\to
\mathcal I_C(6)\to0$, whose positive projective-space cohomology
vanishes \citep[Lemma 30.8.1 and complete proof]{frPrimitiveProjectiveCohomology}. Thus the independent sextic coefficients give a generic
section of $\mathcal O_C(6)$, a degree-$24$ line bundle.

Choose an origin on the constant elliptic curve over $k_0$.
The space of ordered $24$-tuples with sum equal to the class of
$\mathcal O_C(6)$ is isomorphic to $C_K^{23}$ by eliminating the
last point. Removing all diagonals gives a nonempty irreducible
open. Its quotient by the free permutation action is the open of
reduced divisors in $|\mathcal O_C(6)|$; consequently the function
field extension is Galois with group $S_{24}$. This remains true
in every characteristic: the free action of the constant finite
étale group gives an étale torsor on this open, irrespective of
whether $p$ divides $24!$. Adding the independent coefficients in
the kernel of the sextic restriction map preserves this generic
permutation action. Every permutation extends to the common
algebraic closure used for $k$.

For distinct $a,b$, the projection of the fixed-sum ordered-tuple
space to the $(a,b)$ factors is dominant: one may freely choose
those two points and another $21$ points, and solve for the last.
Hence $r_a-r_b$ is nontorsion. All torsion points of the constant
elliptic curve are defined over $k_0$, while these two coordinates
are generic over $K$.

The Picard group of the first blown-up quadric is generated by its
two ruling classes and its exceptional curves. Exactly the $24$
exceptional curves above the $r_a$ meet $C$, and their restrictions
are $\mathcal O_C(r_a)$. All other exceptional curves are disjoint
from $C$. The second component was not blown up at the $C_{12}$
nodes; its blowup centres lie on the other conic away from the
triple points, so its exceptional curves are disjoint from $C$.
Its two ruling restrictions, and those of the first quadric, are
line bundles defined over $k_0$.

Suppose (NP.9) failed. Clearing rational coefficients and then
killing the torsion difference in the integral Picard group gives
an equality
\[
 N[q_j-q_1]=[L_0]+\sum_{a=1}^{24} n_a[r_a]
 \quad\text{in }\operatorname{Pic}(C)(k),
\]
where $N>0$, all $n_a$ are integers, and $L_0$ is defined over
$k_0$. Apply a transposition $(a\ b)$ in the generic permutation
group, fixing $K$ and the sextic coefficients. Subtraction gives
$(n_a-n_b)[r_a-r_b]=0$. Nontorsion forces all $n_a$ to be one
integer $n$. The displayed right hand side is then
$[L_0]+n[\mathcal O_C(6)]$, a degree-zero line bundle defined over
$k_0$. Every point of $\operatorname{Pic}^0(C)(k_0)$ is torsion,
since it is defined over a finite field. This contradicts the
nontorsion of $q_j-q_1$ proved in NP.1--NP.8.
\hfill$\square$
